\label{chap:p3-83-27-menos_un_doceavo}

\section{Domaine Hensel--Witt, applications de Teichmüller et cible rationnelle \(-1/12\)}

Le NTHW est le paquet de données et d’applications
\[
\NTHW=
\bigl(
\mathscr R_\pi^{\rm micro},
\Gnine,
\Cdoce,
K,
\mathcal W_{12\to24},
\mathcal I_{6\to8}
\bigr),
\]
avec les compatibilités suivantes :
\begin{enumerate}
  \item les régions de Hensel partagent un mot quotient et conservent
  des mémoires distinctes ;
  \item \(\Gnine\) transporte l’orientation et produit une monodromie ;
  \item \(\Cdoce\) organise l’observable signé et le sceau \(K\) ;
  \item la lecture de Witt transforme les supports du même bord en hexades et en
  étoiles ;
  \item l’extension \(W_{12}\to W_{24}\) produit l’interface
  hexade--octade ;
  \item les lecteurs archimédien, exceptionnel et dimensionnel reçoivent
  des invariants du même état.
\end{enumerate}

\begin{figure}[!htbp]
\centering
\resizebox{\textwidth}{!}{\begin{tikzpicture}[font=\small,node distance=1.2cm and 1cm,
  nodebox/.style={draw,rounded corners=2mm,align=center,minimum width=3cm,minimum height=1.05cm},
  arr/.style={-{Stealth[length=2.4mm]},thick}]
  \node[nodebox,draw=hmtblue,fill=hmtlightblue] (hensel) at (-5,2.3) {cylindres henséliens\\cocycle de retenue et \(\MonNine\)};
  \node[nodebox,draw=hmtblue,fill=hmtlightblue] (regions) at (-5,-.1) {cinq régions de \(\pi\)\\régions de \(e,\varphi\)};
  \node[nodebox,draw=hmtgold,fill=hmtlightgold,minimum width=3.8cm,minimum height=2.05cm,font=\bfseries] (core) at (0,1.1) {\(\mathscr T_{\rm HW}\)\\Correspondances\\Hensel--Witt};
  \node[nodebox,draw=hmtviolet,fill=hmtlightviolet] (c12) at (0,-1.8) {\(\Cdoce\), \(\Usgn\), \(K\)\\différences \(D_3,D_4,Q\)};
  \node[nodebox,draw=hmtgreen,fill=hmtlightgreen] (witt) at (5,2.3) {Witt \(W_{12}\)\\supports et \(M_{12}\)};
  \node[nodebox,draw=hmtgreen,fill=hmtlightgreen] (oct) at (5,-.1) {inclusion d’incidence \(6\to8\)\\canaux \(90/120\), normalisation \(-1/12\)};
  \node[nodebox,draw=hmtred,fill=hmtlightred] (readers) at (5,-2.5) {lecteurs coordonnés\\valeur et symétrie};
  \draw[arr,hmtblue] (hensel)--(core);
  \draw[arr,hmtblue] (regions)--(core);
  \draw[arr,hmtgold] (core)--(c12);
  \draw[arr,hmtgreen] (core)--(witt);
  \draw[arr,hmtgreen] (core)--(oct);
  \draw[arr,hmtred] (c12)--(readers);
  \draw[arr,hmtred] (witt.east)--++(1.9,0)|-(readers.east);
  \draw[arr,hmtred] (oct)--(readers);
  \node[align=center,hmtgray,font=\footnotesize] at (0,-3.45) {Compatibilité des cylindres, de la dodécaphase et des incidences sur un même état enrichi.};
\end{tikzpicture}
}
\caption{Le NTHW comme interface typée. La tâche mathématique centrale consiste à faire
commuter le diagramme, et non à accumuler des coïncidences autour d’un nom.}
\label{p3-83-c27-s1:fig:nthw}
\end{figure}

Le NTHW donne un nom formel au lieu auparavant décrit comme « pierre de
Rosette ». La substitution est conceptuelle : une pierre traduit par ressemblance ; une
transduction précise l’entrée, la sortie, la donnée conservée et la perte
d’information.

\section{Application de Teichmüller du code ternaire à l’anneau de Witt}

La frontière réversible d’APP sélectionne, avec des jauges déclarées, une matrice
\(A_W\) sur \(\FF_3\) qui satisfait
\[
A_WA_W^T=2I\pmod3.
\]
Le code graphe
\[
C_W=\{(q,qA_W):q\in\FF_3^6\}
\]
est le code de Golay ternaire étendu. Son énumérateur de poids est
\[
1+264y^6+440y^9+24y^{12}.
\]
Les \(264\) mots de poids six, identifiés au signe près, donnent \(132\)
supports. Chaque sous-ensemble de cinq points appartient à une unique hexade :
\[
W_{12}=S(5,6,12).
\]

Le champ de \(495\) involutions étoile engendre un groupe d’ordre
\(95\,040\), reconnu comme \(M_{12}\). Une étoile isolée n’engendre pas le
groupe. Le repère HMT ordonné possède un stabilisateur trivial et fixe donc une
jauge dans l’action.

\section{Incidence hexade--octade et canaux modulaires \(90/120\)}

Soit \(H\) une hexade et \(O_H\) une octade qui la contient. En marquant un point
et une paire de points de l’hexade, on obtient
\[
F_6(H)=H\times\binom H2,
\qquad
|F_6|=6\binom62=90.
\]
Si le point marqué peut parcourir l’octade :
\[
F_8(H)=O_H\times\binom H2,
\qquad
|F_8|=8\binom62=120.
\]
L’inclusion \(H\hookrightarrow O_H\) induit \(F_6\hookrightarrow F_8\).
Son défaut normalisé est
\[
\boxed{
\frac{90-120}{24\binom62}
=\frac{6-8}{24}
=-\frac1{12}.}
\]
Ce théorème combinatoire précède la lecture angulaire de Barbero. Les
deux objets partagent l’interface \(6\to8\), mais ne sont pas identiques.

\section{4 routes compatibles vers \texorpdfstring{\(-1/12\)}{-1/12}}

La construction contient quatre réalisations qui doivent être exposées ensemble et non
comptées comme quatre preuves indépendantes d’une même provenance.

\subsection{Réalisation angulaire par l’incrément élémentaire dodécaphasique}

Une dent de \(\Cdoce\) mesure \(30^\circ\). Alors
\[
\frac{30}{120}-\frac{30}{90}
=\frac14-\frac13
=-\frac1{12}.
\]
C’est l’ombre angulaire du motif \(90/120\).

\subsection{Pseudo-inverse orientée de l’opérateur de Gram sur le sous-espace survivant}

Dans le saut certifié \(60\to81\), deux fibres ont pour tailles \(0\) et
\(12\). L’opérateur de Gram est
\[
K_{60,81}=\operatorname{diag}(0,12).
\]
Sur le sous-espace vivant, sa pseudo-inverse vaut \(1/12\) ; avec l’orientation
négative de bord,
\[
E_{\rm surv}=-K_{60,81}^{+}
\]
produit \(-1/12\).

\subsection{Extraction d’échelle par 2e différence triadique}

La réalisation par un extracteur d’échelle triadique est la construction par
seconde différence développée intégralement dans
\cref{sec:c27-segunda-diferencia-triadica}. Cette première apparition fixe sa
fonction dans l’inventaire des réalisations ; la résidence indiquée conserve
les définitions de \(T_L\), \(a(L)\), \(A_1(L)\), \(C(L)\), l’indépendance
du résidu à l’égard de \(L\) et le statut exact de la normalisation.

\subsection{Réalisation modulaire par le quotient \(\eta(\tau)/\eta(3\tau)\) et la fonction \(t_3\)}

La fonction êta satisfait
\[
\frac{\eta(\tau)}{\eta(3\tau)}
=q^{-1/12}\prod_{n\ge1}(1+q^n+q^{2n})^{-1}.
\]
L’exposant provient de \(1/24-3/24=-1/12\). En prenant douze copies :
\[
t_3(q)=\left(\frac{\eta(\tau)}{\eta(3\tau)}\right)^{12}
=q^{-1}-12+54q-76q^2-243q^3+\cdots.
\]
Le coefficient \(54\) coïncide avec le défaut APP. Le corridor
\[
j=\frac{(t_3+27)(t_3+243)^3}{t_3^3}
\]
se relie ensuite à la branche modulaire classique.

\begin{remark}[Ce qui a été démontré dans le corridor]
Le défaut APP \(54\), l’exposant êta \(-1/12\), le défaut
hexade--octade et les identités modulaires une fois \(t_3\) fixé sont exacts.
L’affirmation forte est qu’ils soient tous des images naturelles d’un unique opérateur
TPK. Cette naturalité doit être démontrée au moyen du diagramme du NTHW ; elle ne se
déduit pas de la seule répétition des mêmes entiers.
\end{remark}


\section{5 réalisations algébriques de \texorpdfstring{\(-1/12\)}{-1/12} et séparation de leurs domaines}

Outre le défaut normalisé de l’inclusion hexade--octade, apparaissent
cinq réalisations opératorielles du même rationnel. Leurs domaines sont
précisés séparément ; l’égalité des valeurs n’identifie pas les
opérateurs sans un entrelaceur supplémentaire.

\subsection{Différence de périodicités dodécaphasiques}

Une dent de \(\Cdoce\) mesure \(30^\circ\). Alors
\[
\frac{30}{120}-\frac{30}{90}
=\frac14-\frac13
=-\frac1{12}.
\]
C’est la différence orientée entre les périodes de \(90^\circ\) et
\(120^\circ\) dans cette réalisation.

\subsection{Pseudo-inverse de l’opérateur de Gram sur la composante vivante}

Dans le saut exact \(60\to81\), deux fibres ont pour tailles \(0\) et
\(12\). L’opérateur de Gram est
\[
K_{60,81}=\operatorname{diag}(0,12).
\]
Sur la composante non nulle de son image, la pseudo-inverse vaut \(1/12\) ; avec l’orientation
négative de bord,
\[
E_{\rm surv}=-K_{60,81}^{+}
\]
produit \(-1/12\).

\subsection{2e différence entre échelles triadiques}
\label{sec:c27-segunda-diferencia-triadica}

Définissons
\[
T_L=\sum_{n=1}^{L-1}n(L-n)=\frac{L(L^2-1)}6,
\qquad
a(L)=-\frac{T_L}{L^2}=-\frac L6+\frac1{6L}.
\]
La première soustraction d’échelle est
\[
A_1(L)=a(L)-\frac{a(3L)}3=\frac4{27L},
\]
et la seconde
\[
C(L)=LA_1(L)-\frac{3L}{3}A_1(3L)=\frac8{81}.
\]
Le résidu \(8/81\) est indépendant de \(L\). Pour obtenir
\(-1/12\), on applique la normalisation déclarée
\[
R(L)=-\frac{27}{32}C(L)=-\frac1{12}.
\]
Le calcul dérive \(8/81\) ; il ne dérive pas à lui seul le facteur \(-27/32\).

\subsection{Cadena modular}

La fonction êta satisfait
\[
\frac{\eta(\tau)}{\eta(3\tau)}
=q^{-1/12}\prod_{n\ge1}(1+q^n+q^{2n})^{-1}.
\]
L’exposant provient de \(1/24-3/24=-1/12\). En prenant douze copies :
\[
t_3(q)=\left(\frac{\eta(\tau)}{\eta(3\tau)}\right)^{12}
=q^{-1}-12+54q-76q^2-243q^3+\cdots.
\]
Le coefficient \(54\) coïncide avec le défaut APP\@. L’identité rationnelle
\[
j=\frac{(t_3+27)(t_3+243)^3}{t_3^3}
\]
établit la connexion avec la branche modulaire classique.

\Needspace{32\baselineskip}
\subsection{Défaut des projecteurs de différence cyclique}

Soit \(S\) le décalage cyclique sur \(\QQ^{12}\) et définissons
\[
 D_r=I-S^r.
\]
Le noyau de \(D_r\) est formé des vecteurs constants sur chaque orbite
de \(S^r\) ; par conséquent,
\[
 \dim\ker D_r=\gcd(12,r).
\]
Soient \(\Pi_{\ker D_3}\) et \(\Pi_{\ker D_4}\) les projecteurs rationnels sur
\(\ker D_3\) et \(\ker D_4\), et soit
\(\tau(T)=\operatorname{Tr}(T)/12\) la trace normalisée. Alors
\[
 \boxed{
 \tau(\Pi_{\ker D_3}-\Pi_{\ker D_4})
 =\frac{\operatorname{rank}\Pi_{\ker D_3}
       -\operatorname{rank}\Pi_{\ker D_4}}{12}
 =\frac{3-4}{12}
 =-\frac1{12}.}
\]
Cette réalisation exprime le rationnel comme défaut transversal des pas
trois et quatre dans le cycle à douze phases.

\begin{remark}[Domaines des réalisations]
Le défaut APP \(54\), l’exposant êta \(-1/12\), le défaut
combinatoire du relèvement hexade--octade et les
identités modulaires une fois \(t_3\) fixé sont exacts. La valeur
de la pseudo-inverse, l’extracteur normalisé et le défaut des projecteurs sont également exacts.
Chaque égalité conserve son domaine et son application correspondante.
\end{remark}
